\section{Related Work}\label{sec:related_work}

\subsection{The Collapse of K-Anonymity}
Traditional anonymization relies on the removal or suppression of direct identifiers such as IMSI numbers, phone numbers, and subscriber names. However, research has demonstrated that these techniques are fundamentally insufficient when spatial behavior remains stable \cite{duan2025sure}. The SURE framework shows that even with k-anonymity applied to shared datasets, the underlying spatio-temporal metadata BS associations, traffic patterns, and temporal routines persists unchanged, enabling re-identification with AUC scores exceeding 0.9 \cite{duan2025sure}.

A critical finding from the SURE study is that user demographics significantly impact re-identification susceptibility. Male users and those aged 45 and above exhibit the highest re-identification risks (AUC scores of 0.95 for males and above 0.93 for users over 45), attributable to their more consistent and predictable logistical routines \cite{duan2025sure}. Conversely, users aged 18--24 show marginally lower risk, likely due to more variable mobility patterns. These findings demonstrate that the ``Top 2 BS'' spatial anchor, where the majority of a user's traffic concentrates on just two base stations, provides a near-deterministic fingerprint of home and workplace locations that survives any identifier-level anonymization \cite{duan2025sure}.

Furthermore, MoCo \cite{duan2025moco} extends the threat beyond individual re-identification to social graph reconstruction. By detecting co-traveling patterns with an F1-score of 87.13\%, MoCo can map a target's ``inner circle''the set of users who regularly share transportation routes using only cellular signaling traces \cite{duan2025moco}. This compound threat renders k-anonymity not only insufficient for identity protection but also incapable of shielding social associations.

\subsection{Alternative SSL Defenses}
Several recent works have proposed defensive techniques against adversarial attacks on self-supervised learning (SSL) models. ASTRA (Adversarial Self-Supervised Training with Adaptive-Attacks) \cite{chhipa2025astra} introduces an adaptive adversarial training procedure that generates perturbations during the pre-training phase, improving robustness against adversarial inputs. Similarly, RoCL (Robust Contrastive Learning) \cite{kim2020rocl} incorporates adversarial examples directly into the contrastive learning objective,
training the encoder to produce consistent representations even under input perturbations. Additionally, DIFT \cite{zhu2026dift} proposes density-based identification and fine-tuning to counter data poisoning attacks in contrastive learning.

The success of these architectures yields a critical insight for privacy and tracking evasion: SSL-based models are highly adept at adapting to sub-optimal, 
noisy, or perturbed data conditions. Crucially, as long as a tracking model is trained with potential data variations in mind, it can maintain high efficacy even without prior knowledge 
of the exact noise distribution. Consequently, attempts to evade tracking by injecting noise or obfuscating spatio-temporal metadata are fundamentally flawed. 
An adversary operating at the BSS/OSS database level  can simply leverage these robust SSL frameworks to train tracking models (such as SURE or MoCo) that easily bypass data-level perturbations.
This capability indicates that data-obfuscation defenses may not guarantee privacy against adaptive tracking models. Therefore, the most reliable method to avoid being tracked is
to completely circumvent the data collection infrastructure. This motivates our proposed approach: deploying a parallel communication network that inherently avoids generating the centralized spatio-temporal 
metadata exploited by these models in the first place.
